\section{Konsep Dasar Fungsi Hash}

Bab-bab sebelum ini membahas kerahasiaan: bagaimana menyembunyikan isi pesan dari
mata yang tidak berhak. Namun kerahasiaan bukan satu-satunya layanan keamanan yang
dibutuhkan. Sering kali isi pesan justru boleh diketahui siapa saja, asalkan
dapat dipastikan bahwa isi itu \emph{tidak berubah} dan bahwa ia datang dari
pihak yang benar. Inilah masalah \textbf{integritas data}, dan perkakas utamanya
adalah \textbf{fungsi hash}.

Fungsi hash adalah algoritma yang memetakan data berukuran apa pun --- satu huruf,
satu berkas berukuran gigabita, atau seluruh basis data --- menjadi keluaran
berukuran tetap (\textit{fixed size}) yang disebut \textit{hash value},
\textit{message digest}, atau \textit{fingerprint}. Berbeda dengan enkripsi yang
bersifat dua arah (ada enkripsi, ada dekripsi), fungsi hash adalah \textbf{fungsi
satu arah} (\textit{one-way function}). Menghitung hash dari sebuah pesan selalu
mudah; mengembalikan pesan asli dari nilai hash-nya, secara komputasi, tidak dapat
dilakukan.

% Alur fungsi hash: pesan berukuran bebas dipetakan menjadi digest berukuran tetap.
% Dirujuk dari section-01.tex dengan % Alur fungsi hash: pesan berukuran bebas dipetakan menjadi digest berukuran tetap.
% Dirujuk dari section-01.tex dengan % Alur fungsi hash: pesan berukuran bebas dipetakan menjadi digest berukuran tetap.
% Dirujuk dari section-01.tex dengan \input{figures/alur-fungsi-hash}.
\begin{figure}[htbp]
    \centering
    \begin{tikzpicture}[
        kotak/.style={draw, rounded corners, align=center, font=\small, inner sep=6pt},
        panah/.style={-{Latex[length=2.4mm]}, thick},
        node distance=1.4cm,
    ]
        \node[kotak, fill=blue!6, minimum width=3.4cm, minimum height=1.5cm] (m)
            {Pesan $M$\\[2pt]\scriptsize panjang bebas};
        \node[kotak, fill=orange!12, right=of m, minimum width=3.2cm, minimum height=1.9cm] (h)
            {Fungsi hash\\[2pt]$h = H(M)$};
        \node[kotak, fill=green!10, right=of h, minimum width=3.2cm, minimum height=1.5cm] (d)
            {Digest $h$\\[2pt]\scriptsize panjang tetap};
        \draw[panah] (m) -- (h);
        \draw[panah] (h) -- (d);
    \end{tikzpicture}
    \caption{Alur kerja fungsi hash: pesan berukuran bebas dipetakan menjadi digest berukuran tetap}
    \label{fig:alur-fungsi-hash}
\end{figure}
.
\begin{figure}[htbp]
    \centering
    \begin{tikzpicture}[
        kotak/.style={draw, rounded corners, align=center, font=\small, inner sep=6pt},
        panah/.style={-{Latex[length=2.4mm]}, thick},
        node distance=1.4cm,
    ]
        \node[kotak, fill=blue!6, minimum width=3.4cm, minimum height=1.5cm] (m)
            {Pesan $M$\\[2pt]\scriptsize panjang bebas};
        \node[kotak, fill=orange!12, right=of m, minimum width=3.2cm, minimum height=1.9cm] (h)
            {Fungsi hash\\[2pt]$h = H(M)$};
        \node[kotak, fill=green!10, right=of h, minimum width=3.2cm, minimum height=1.5cm] (d)
            {Digest $h$\\[2pt]\scriptsize panjang tetap};
        \draw[panah] (m) -- (h);
        \draw[panah] (h) -- (d);
    \end{tikzpicture}
    \caption{Alur kerja fungsi hash: pesan berukuran bebas dipetakan menjadi digest berukuran tetap}
    \label{fig:alur-fungsi-hash}
\end{figure}
.
\begin{figure}[htbp]
    \centering
    \begin{tikzpicture}[
        kotak/.style={draw, rounded corners, align=center, font=\small, inner sep=6pt},
        panah/.style={-{Latex[length=2.4mm]}, thick},
        node distance=1.4cm,
    ]
        \node[kotak, fill=blue!6, minimum width=3.4cm, minimum height=1.5cm] (m)
            {Pesan $M$\\[2pt]\scriptsize panjang bebas};
        \node[kotak, fill=orange!12, right=of m, minimum width=3.2cm, minimum height=1.9cm] (h)
            {Fungsi hash\\[2pt]$h = H(M)$};
        \node[kotak, fill=green!10, right=of h, minimum width=3.2cm, minimum height=1.5cm] (d)
            {Digest $h$\\[2pt]\scriptsize panjang tetap};
        \draw[panah] (m) -- (h);
        \draw[panah] (h) -- (d);
    \end{tikzpicture}
    \caption{Alur kerja fungsi hash: pesan berukuran bebas dipetakan menjadi digest berukuran tetap}
    \label{fig:alur-fungsi-hash}
\end{figure}


\subsection{Definisi dan Notasi}

Secara formal, fungsi hash kriptografis adalah pemetaan

\begin{equation}
H : \{0, 1\}^{*} \longrightarrow \{0, 1\}^{n},
\label{eq:hash-definisi}
\end{equation}

yang menerima untai bit dengan panjang sembarang (termasuk panjang nol) dan
mengembalikan untai bit tetap sepanjang $n$. Nilai $n$ disebut panjang digest, dan
besarnya berbeda-beda antara algoritma: MD5 memakai $n = 128$, SHA-1 memakai
$n = 160$, sedangkan keluarga SHA-2 dan SHA-3 menyediakan $n = 224$, $256$, $384$,
dan $512$.

Perhatikan bahwa domain $H$ jauh lebih besar daripada kodomainnya. Ada tak
berhingga banyak pesan yang mungkin dan hanya $2^{n}$ nilai hash yang mungkin,
sehingga pasti ada pasangan pesan berbeda dengan hash yang sama. Kenyataan itu
bukan cacat perancangan, melainkan hukum matematika yang tidak dapat dielakkan
oleh algoritma apa pun. Yang dapat dirancang hanyalah \emph{seberapa sukar}
pasangan semacam itu ditemukan, dan di situlah letak seluruh persoalan.

Satu gagasan lagi diperlukan sebelum properti keamanan dapat dirumuskan. Sebagian
besar fungsi hash praktis tidak memproses pesan sepanjang apa pun dalam satu
langkah. Karena itu fungsi hash modern dibangun dari sebuah \textbf{fungsi
kompresi} (\textit{compression function}) $f$ yang menerima blok pesan berukuran
tetap dan mengubah keadaan internal berukuran tetap pula. Properti keamanan yang
dirumuskan untuk $H$ sesungguhnya bertumpu pada properti $f$, dan hal ini akan
dibahas lebih jauh pada Subbagian~\ref{subsec:padding-hash} dan
Subbagian~\ref{subsec:konstruksi-hash}.

\subsection{Karakteristik Fungsi Hash Kriptografis}

Sebuah fungsi hash dikatakan layak dipakai untuk keperluan kriptografi jika
memenuhi properti berikut:

\begin{enumerate}
    \item \textbf{Efisiensi}: Perhitungan hash untuk input apa pun harus dapat
          dilakukan dengan cepat.
    \item \textbf{Pre-image Resistance}: Diberikan nilai hash $h$, sangat sulit
          untuk menemukan pesan $m$ sedemikian sehingga $H(m) = h$.
    \item \textbf{Second Pre-image Resistance}: Diberikan pesan $m_1$, sangat
          sulit untuk menemukan pesan $m_2 \neq m_1$ sedemikian sehingga
          $H(m_1) = H(m_2)$.
    \item \textbf{Collision Resistance}: Sangat sulit untuk menemukan dua pesan
          berbeda $m_1$ dan $m_2$ sedemikian sehingga $H(m_1) = H(m_2)$.
\end{enumerate}

Ketiga properti keamanan itu tampak serupa, tetapi pertanyaannya berbeda dan
biayanya pun berbeda. Perbedaan itu paling jelas bila dituliskan sebagai
persoalan pencarian:

\begin{align}
\text{\textit{Pre-image}:}   \quad & \text{diberikan } h, \text{ cari } m
   \text{ dengan } H(m) = h ,
   \label{eq:preimage} \\[2pt]
\text{\textit{Second pre-image}:} \quad & \text{diberikan } m_1, \text{ cari }
   m_2 \neq m_1 \text{ dengan } H(m_2) = H(m_1) ,
   \label{eq:second-preimage} \\[2pt]
\text{\textit{Kolisi}:} \quad & \text{cari sebarang pasangan } m_1 \neq m_2
   \text{ dengan } H(m_1) = H(m_2) .
   \label{eq:kolisi}
\end{align}

Pada persoalan pre-image dan second pre-image, sasarannya sudah ditetapkan lebih
dahulu: penyerang harus menemukan pesan yang hash-nya tepat sama dengan sebuah
nilai tertentu. Peluang setiap tebakan adalah $1/2^{n}$, sehingga dibutuhkan
sekitar $2^{n}$ percobaan. Pada persoalan kolisi, sasarannya tidak ditetapkan:
penyerang bebas memilih \emph{apa saja} yang penting dua pesannya berhash sama.
Kebebasan itu mengubah biayanya secara drastis, dan penyebabnya adalah paradoks
ulang tahun.

Tabel~\ref{tab:ketahanan-hash} merangkum ketiga properti itu beserta biaya
pencarian yang diharapkan dan kegunaannya masing-masing.

\begin{table}[htbp]
    \centering
    \caption{Tiga properti keamanan fungsi hash dan biaya penyerangannya}
    \label{tab:ketahanan-hash}
    \begin{tabularx}{\linewidth}{@{}l
                                   >{\raggedright\arraybackslash}X
                                   >{\raggedright\arraybackslash}X@{}}
    \toprule
    \textbf{Properti} & \textbf{Pertanyaan penyerang} &
    \textbf{Biaya} \\
    \midrule
    Pre-image & Diberikan digest $h$, temukan pesan apa pun yang berhash $h$ &
    $2^{n}$ \\
    \addlinespace
    Second pre-image & Diberikan pesan $m_1$, temukan pesan lain berhash sama &
    $2^{n}$ \\
    \addlinespace
    Kolisi & Temukan sebarang dua pesan berbeda yang berhash sama &
    $2^{n/2}$ \\
    \bottomrule
    \end{tabularx}
\end{table}

Pembedaan itu bukan perbedaan kecil. Untuk SHA-256 ($n = 256$), menemukan
pre-image menuntut sekitar $2^{256}$ percobaan --- jauh di luar jangkauan
perangkat apa pun yang dapat dibayangkan. Namun menemukan kolisi hanya menuntut
sekitar $2^{128}$ percobaan. Angka $2^{128}$ itu tetap sangat besar, tetapi ia
berbeda dari $2^{256}$ bukan sekadar beberapa kali, melainkan sebesar
$2^{128}$ kali. Inilah sebabnya sebuah fungsi hash dengan digest 256 bit hanya
menawarkan keamanan kolisi setara 128 bit, dan itulah tingkat keamanan yang
sesungguhnya harus dibandingkan dengan panjang kunci algoritma lain.

\subsection{Serangan Birthday dan Batas Nyata Ketahanan Kolisi}
\label{subsec:birthday}

Ketahanan kolisi tidak sebesar yang disarankan ukuran digest. Karena
\textit{birthday paradox}, menemukan dua pesan berhash sama hanya membutuhkan
sekitar $\sqrt{2^n} = 2^{n/2}$ percobaan acak, bukan $2^n$
\citep{stallings2017,menezes1996}.

Paradoks ulang tahun adalah nama bagi kenyataan bahwa dalam himpunan yang cukup
besar, peluang dua anggota berulang jauh lebih tinggi daripada dugaan. Bila
sebuah ruangan berisi $k$ orang dan setiap ulang tahun dianggap tersebar merata
pada 365 hari, maka peluang ada dua orang berulang tahun sama bukanlah
$k/365$ --- melainkan jauh lebih besar. Sebabnya, yang dihitung bukan peluang
satu orang tertentu bertabrakan dengan orang tertentu lain, melainkan peluang
\emph{ada sepasang} di antara seluruh anggota yang bertabrakan. Jumlah pasangan
tumbuh kuadratik terhadap $k$, yaitu $k(k-1)/2$, sehingga peluangnya membesar
jauh lebih cepat daripada dugaan linier.

Hal yang sama berlaku pada fungsi hash. Untuk $k$ pesan acak, peluang terjadinya
kolisi adalah

\begin{equation}
P(\text{kolisi}) \approx 1 - e^{-k(k-1)/(2\cdot 2^n)}
   \approx 1 - e^{-k^2/2^{\,n+1}}
   \approx \frac{k^2}{2^{\,n+1}} ,
\label{eq:birthday-peluang}
\end{equation}

dengan pendekatan terakhir berlaku selama peluangnya masih kecil. Bentuk kedua
itulah yang menarik untuk dibalik: berapa banyak pesan harus dicoba agar peluang
kolisi mencapai setengah? Substitusi $P = 1/2$ memberi

\begin{equation}
k \approx \sqrt{2 \ln 2 \cdot 2^{\,n}}
  = \sqrt{2 \ln 2} \cdot 2^{\,n/2}
  \approx 1{,}177 \cdot 2^{\,n/2} .
\label{eq:birthday-50}
\end{equation}

Faktor $\sqrt{2\ln 2} = 1{,}1774\ldots$ berasal dari pemilihan peluang $50\%$; bila
yang dikehendaki peluang $1\%$, faktornya tinggal $0{,}2$. Yang menentukan bukan
faktornya, melainkan pangkat $2^{\,n/2}$: biaya selalu sebanding dengan akar
kuadrat ukuran ruang hash.

Sifat itu dapat dilihat sekaligus. Bila banyaknya pesan dinyatakan relatif
terhadap $2^{\,n/2}$, maka Persamaan~\eqref{eq:birthday-peluang} hanya bergantung
pada nisbah itu dan tidak lagi pada $n$. Karena itu kurva pada
Gambar~\ref{fig:batas-birthday} berlaku untuk \emph{semua} panjang digest:
mengganti MD5 dengan SHA-256 tidak mengubah bentuk kurvanya, melainkan hanya
menggeser letaknya pada sumbu mendatar sejauh $2^{64}$. Inilah gambaran visual
dari ``keamanan efektif $n/2$ bit''.

% Batas ulang tahun: peluang kolisi sebagai fungsi banyaknya pesan.
% Sumbu-x dinyatakan relatif terhadap 2^{n/2}, sehingga kurvanya berlaku untuk
% semua panjang digest n. Koordinat dihitung dengan Python dari
% P(u) = 1 - exp(-2^{2u-1}), u = log2(k / 2^{n/2}).
% Dirujuk dari section-01.tex dengan % Batas ulang tahun: peluang kolisi sebagai fungsi banyaknya pesan.
% Sumbu-x dinyatakan relatif terhadap 2^{n/2}, sehingga kurvanya berlaku untuk
% semua panjang digest n. Koordinat dihitung dengan Python dari
% P(u) = 1 - exp(-2^{2u-1}), u = log2(k / 2^{n/2}).
% Dirujuk dari section-01.tex dengan % Batas ulang tahun: peluang kolisi sebagai fungsi banyaknya pesan.
% Sumbu-x dinyatakan relatif terhadap 2^{n/2}, sehingga kurvanya berlaku untuk
% semua panjang digest n. Koordinat dihitung dengan Python dari
% P(u) = 1 - exp(-2^{2u-1}), u = log2(k / 2^{n/2}).
% Dirujuk dari section-01.tex dengan \input{figures/batas-birthday}.
\begin{figure}[htbp]
    \centering
    \begin{tikzpicture}[
        sumbu/.style={-{Latex[length=2.2mm]}, thick},
        kisi/.style={gray!35, very thin},
        tanda/.style={red!70!black, dashed, thick},
    ]
        % kisi mendatar
        \foreach \y in {0.8, 1.6, 2.4} {
            \draw[kisi] (-4.6, \y) -- (2.5, \y);
        }
        % sumbu
        \draw[sumbu] (-4.6, 0) -- (2.7, 0) node[right, font=\small] {$k / 2^{\,n/2}$};
        \draw[sumbu] (0, 0) -- (0, 3.45) node[above, font=\small] {$P(\text{kolisi})$};

        % penanda skala pada sumbu-x
        \foreach \x/\lab in {-2.88/{$2^{-4}$}, -1.44/{$2^{-2}$}, 0/{$2^{0}$}, 1.44/{$2^{2}$}} {
            \draw[kisi] (\x, 0) -- (\x, -0.08);
            \node[below, font=\scriptsize] at (\x, -0.08) {\lab};
        }
        % penanda skala pada sumbu-y
        \foreach \y/\lab in {0.8/{$0{,}25$}, 1.6/{$0{,}50$}, 2.4/{$0{,}75$}} {
            \draw[kisi] (0, \y) -- (-0.08, \y);
            \node[left, font=\scriptsize] at (-0.08, \y) {\lab};
        }

        % kurva peluang kolisi
        \draw[very thick, blue!65!black]
            plot coordinates {
            (-5.760,0.0000) (-5.580,0.0000) (-5.400,0.0000) (-5.220,0.0001)
            (-5.040,0.0001) (-4.860,0.0001) (-4.680,0.0002) (-4.500,0.0003)
            (-4.320,0.0004) (-4.140,0.0006) (-3.960,0.0008) (-3.780,0.0011)
            (-3.600,0.0016) (-3.420,0.0022) (-3.240,0.0031) (-3.060,0.0044)
            (-2.880,0.0062) (-2.700,0.0088) (-2.520,0.0125) (-2.340,0.0176)
            (-2.160,0.0249) (-1.980,0.0352) (-1.800,0.0496) (-1.620,0.0699)
            (-1.440,0.0985) (-1.260,0.1383) (-1.080,0.1939) (-0.900,0.2707)
            (-0.720,0.3760) (-0.540,0.5185) (-0.360,0.7078) (-0.180,0.9530)
            (0.000,1.2591) (0.180,1.6222) (0.360,2.0228) (0.540,2.4220)
            (0.720,2.7669) (0.900,3.0109) (1.080,3.1414) (1.260,3.1888)
            (1.440,3.1989) (1.620,3.2000) (1.800,3.2000) (1.980,3.2000)
            (2.160,3.2000) (2.340,3.2000) (2.520,3.2000) (2.700,3.2000)
            (2.880,3.2000)
            };

        % penanda $k = 2^{n/2}$
        \draw[tanda] (0, 0) -- (0, 1.2591);
        \node[fill=white, draw=red!70!black, font=\scriptsize, align=left,
              inner sep=2.5pt, right] at (0.06, 0.62)
            {$k = 2^{\,n/2}$:\\ $P \approx 39{,}3\%$};

        % penanda peluang 50%
        \draw[tanda] (0.169, 1.6) -- (1.44, 1.6);
        \node[font=\scriptsize, align=left, anchor=west] at (1.48, 1.42)
            {$P = 50\%$\\ pada\\ $k = 1{,}177 \cdot 2^{\,n/2}$};
        \draw[red!70!black, thick] (0.169, 1.585) -- (0.169, 1.615);
    \end{tikzpicture}
    \caption{Peluang terjadinya kolisi sebagai fungsi banyaknya pesan yang dicoba,
    dinyatakan relatif terhadap $2^{\,n/2}$. Kurva ini berlaku untuk semua panjang
    digest $n$ --- perubahan $n$ hanya menggeser letaknya pada sumbu mendatar}
    \label{fig:batas-birthday}
\end{figure}
.
\begin{figure}[htbp]
    \centering
    \begin{tikzpicture}[
        sumbu/.style={-{Latex[length=2.2mm]}, thick},
        kisi/.style={gray!35, very thin},
        tanda/.style={red!70!black, dashed, thick},
    ]
        % kisi mendatar
        \foreach \y in {0.8, 1.6, 2.4} {
            \draw[kisi] (-4.6, \y) -- (2.5, \y);
        }
        % sumbu
        \draw[sumbu] (-4.6, 0) -- (2.7, 0) node[right, font=\small] {$k / 2^{\,n/2}$};
        \draw[sumbu] (0, 0) -- (0, 3.45) node[above, font=\small] {$P(\text{kolisi})$};

        % penanda skala pada sumbu-x
        \foreach \x/\lab in {-2.88/{$2^{-4}$}, -1.44/{$2^{-2}$}, 0/{$2^{0}$}, 1.44/{$2^{2}$}} {
            \draw[kisi] (\x, 0) -- (\x, -0.08);
            \node[below, font=\scriptsize] at (\x, -0.08) {\lab};
        }
        % penanda skala pada sumbu-y
        \foreach \y/\lab in {0.8/{$0{,}25$}, 1.6/{$0{,}50$}, 2.4/{$0{,}75$}} {
            \draw[kisi] (0, \y) -- (-0.08, \y);
            \node[left, font=\scriptsize] at (-0.08, \y) {\lab};
        }

        % kurva peluang kolisi
        \draw[very thick, blue!65!black]
            plot coordinates {
            (-5.760,0.0000) (-5.580,0.0000) (-5.400,0.0000) (-5.220,0.0001)
            (-5.040,0.0001) (-4.860,0.0001) (-4.680,0.0002) (-4.500,0.0003)
            (-4.320,0.0004) (-4.140,0.0006) (-3.960,0.0008) (-3.780,0.0011)
            (-3.600,0.0016) (-3.420,0.0022) (-3.240,0.0031) (-3.060,0.0044)
            (-2.880,0.0062) (-2.700,0.0088) (-2.520,0.0125) (-2.340,0.0176)
            (-2.160,0.0249) (-1.980,0.0352) (-1.800,0.0496) (-1.620,0.0699)
            (-1.440,0.0985) (-1.260,0.1383) (-1.080,0.1939) (-0.900,0.2707)
            (-0.720,0.3760) (-0.540,0.5185) (-0.360,0.7078) (-0.180,0.9530)
            (0.000,1.2591) (0.180,1.6222) (0.360,2.0228) (0.540,2.4220)
            (0.720,2.7669) (0.900,3.0109) (1.080,3.1414) (1.260,3.1888)
            (1.440,3.1989) (1.620,3.2000) (1.800,3.2000) (1.980,3.2000)
            (2.160,3.2000) (2.340,3.2000) (2.520,3.2000) (2.700,3.2000)
            (2.880,3.2000)
            };

        % penanda $k = 2^{n/2}$
        \draw[tanda] (0, 0) -- (0, 1.2591);
        \node[fill=white, draw=red!70!black, font=\scriptsize, align=left,
              inner sep=2.5pt, right] at (0.06, 0.62)
            {$k = 2^{\,n/2}$:\\ $P \approx 39{,}3\%$};

        % penanda peluang 50%
        \draw[tanda] (0.169, 1.6) -- (1.44, 1.6);
        \node[font=\scriptsize, align=left, anchor=west] at (1.48, 1.42)
            {$P = 50\%$\\ pada\\ $k = 1{,}177 \cdot 2^{\,n/2}$};
        \draw[red!70!black, thick] (0.169, 1.585) -- (0.169, 1.615);
    \end{tikzpicture}
    \caption{Peluang terjadinya kolisi sebagai fungsi banyaknya pesan yang dicoba,
    dinyatakan relatif terhadap $2^{\,n/2}$. Kurva ini berlaku untuk semua panjang
    digest $n$ --- perubahan $n$ hanya menggeser letaknya pada sumbu mendatar}
    \label{fig:batas-birthday}
\end{figure}
.
\begin{figure}[htbp]
    \centering
    \begin{tikzpicture}[
        sumbu/.style={-{Latex[length=2.2mm]}, thick},
        kisi/.style={gray!35, very thin},
        tanda/.style={red!70!black, dashed, thick},
    ]
        % kisi mendatar
        \foreach \y in {0.8, 1.6, 2.4} {
            \draw[kisi] (-4.6, \y) -- (2.5, \y);
        }
        % sumbu
        \draw[sumbu] (-4.6, 0) -- (2.7, 0) node[right, font=\small] {$k / 2^{\,n/2}$};
        \draw[sumbu] (0, 0) -- (0, 3.45) node[above, font=\small] {$P(\text{kolisi})$};

        % penanda skala pada sumbu-x
        \foreach \x/\lab in {-2.88/{$2^{-4}$}, -1.44/{$2^{-2}$}, 0/{$2^{0}$}, 1.44/{$2^{2}$}} {
            \draw[kisi] (\x, 0) -- (\x, -0.08);
            \node[below, font=\scriptsize] at (\x, -0.08) {\lab};
        }
        % penanda skala pada sumbu-y
        \foreach \y/\lab in {0.8/{$0{,}25$}, 1.6/{$0{,}50$}, 2.4/{$0{,}75$}} {
            \draw[kisi] (0, \y) -- (-0.08, \y);
            \node[left, font=\scriptsize] at (-0.08, \y) {\lab};
        }

        % kurva peluang kolisi
        \draw[very thick, blue!65!black]
            plot coordinates {
            (-5.760,0.0000) (-5.580,0.0000) (-5.400,0.0000) (-5.220,0.0001)
            (-5.040,0.0001) (-4.860,0.0001) (-4.680,0.0002) (-4.500,0.0003)
            (-4.320,0.0004) (-4.140,0.0006) (-3.960,0.0008) (-3.780,0.0011)
            (-3.600,0.0016) (-3.420,0.0022) (-3.240,0.0031) (-3.060,0.0044)
            (-2.880,0.0062) (-2.700,0.0088) (-2.520,0.0125) (-2.340,0.0176)
            (-2.160,0.0249) (-1.980,0.0352) (-1.800,0.0496) (-1.620,0.0699)
            (-1.440,0.0985) (-1.260,0.1383) (-1.080,0.1939) (-0.900,0.2707)
            (-0.720,0.3760) (-0.540,0.5185) (-0.360,0.7078) (-0.180,0.9530)
            (0.000,1.2591) (0.180,1.6222) (0.360,2.0228) (0.540,2.4220)
            (0.720,2.7669) (0.900,3.0109) (1.080,3.1414) (1.260,3.1888)
            (1.440,3.1989) (1.620,3.2000) (1.800,3.2000) (1.980,3.2000)
            (2.160,3.2000) (2.340,3.2000) (2.520,3.2000) (2.700,3.2000)
            (2.880,3.2000)
            };

        % penanda $k = 2^{n/2}$
        \draw[tanda] (0, 0) -- (0, 1.2591);
        \node[fill=white, draw=red!70!black, font=\scriptsize, align=left,
              inner sep=2.5pt, right] at (0.06, 0.62)
            {$k = 2^{\,n/2}$:\\ $P \approx 39{,}3\%$};

        % penanda peluang 50%
        \draw[tanda] (0.169, 1.6) -- (1.44, 1.6);
        \node[font=\scriptsize, align=left, anchor=west] at (1.48, 1.42)
            {$P = 50\%$\\ pada\\ $k = 1{,}177 \cdot 2^{\,n/2}$};
        \draw[red!70!black, thick] (0.169, 1.585) -- (0.169, 1.615);
    \end{tikzpicture}
    \caption{Peluang terjadinya kolisi sebagai fungsi banyaknya pesan yang dicoba,
    dinyatakan relatif terhadap $2^{\,n/2}$. Kurva ini berlaku untuk semua panjang
    digest $n$ --- perubahan $n$ hanya menggeser letaknya pada sumbu mendatar}
    \label{fig:batas-birthday}
\end{figure}


Angka konkretnya memperlihatkan akibat pembedaan itu kepada setiap algoritma:

\begin{itemize}
    \item MD5 ($n=128$): $k \approx 1{,}177 \cdot 2^{64} \approx
          2{,}2\times10^{19}$, sebab
          \[ 2^{64} = 18.446.744.073.709.551.616 . \]
    \item SHA-1 ($n=160$): $k \approx 1{,}177 \cdot 2^{80} \approx
          1{,}4\times10^{24}$.
    \item SHA-256 ($n=256$): $k \approx 1{,}177 \cdot 2^{128} \approx
          4{,}0\times10^{38}$.
    \item SHA-512 ($n=512$): $k \approx 1{,}177 \cdot 2^{256} \approx
          1{,}4\times10^{77}$.
\end{itemize}

\important{Ini berarti kekuatan efektif terhadap kolisi hanya $n/2$ bit, bukan
$n$ bit.} SHA-256 setara 128 bit dan masih aman, sedangkan MD5 setara 64 bit dan
sudah dapat ditembus GPU \citep{menezes1996}. Karena itu properti yang menentukan
pilihan hash untuk tanda tangan digital adalah \textit{collision resistance},
bukan pre-image resistance.

Perlu ditegaskan bahwa angka-angka di atas adalah batas \emph{generik}, yaitu yang
berlaku bagi fungsi hash yang berperilaku seperti fungsi acak. Sebuah fungsi hash
yang menyimpang dari perilaku acak dapat diserang jauh lebih murah lagi, dan itulah
yang terjadi pada MD5 dan SHA-1: serangan terhadap keduanya bekerja melalui
kelemahan struktural pada fungsi kompresinya, bukan melalui pencarian
menyeluruh. Serangan-serangan itu dibahas pada Subbagian~\ref{subsec:kriptanalisis-hash}.

\subsection{Panjang Pesan dan Padding: Mengapa Kolisi Menular}
\label{subsec:padding-hash}

Sampai di sini fungsi hash masih dibayangkan sebagai kotak ajaib yang menerima
pesan sepanjang apa pun. Kenyataannya tidak demikian. Fungsi kompresi bekerja
atas blok berukuran tetap --- 512 bit pada MD5, SHA-1, dan SHA-256, serta 1024
bit pada SHA-384 dan SHA-512 --- sehingga pesan harus lebih dahulu dipotong
menjadi blok. Potongan terakhir hampir selalu tidak penuh, dan ketidakpenuhan itu
harus diselesaikan dengan aturan yang tidak dapat menimbulkan kerancuan.

Aturan itu disebut \textbf{padding} dan mengikuti pola yang sama pada seluruh
keluarga SHA-2. Bit tunggal \code{1} disisipkan tepat sesudah pesan, lalu diikuti
sebanyak mungkin bit \code{0} sehingga panjang totalnya bersisa $448$ modulo
$512$, dan akhirnya panjang pesan semula dituliskan sebagai bilangan 64 bit.
Dengan aturan itu, panjang pesan yang diproses selalu kelipatan bulat dari 512
bit, dan satu blok penuh selalu tersisa untuk memuat catatan panjangnya:

\begin{equation}
\text{blok akhir} = \underbrace{m}_{\text{sisa pesan}} \,\|\, \code{1} \,\|\,
\underbrace{\code{0}\cdots\code{0}}_{\text{hingga bit ke-}448}
\,\|\, \underbrace{\mathrm{panjang}(m)\vphantom{\big|}}_{\text{64 bit}} .
\label{eq:padding-sha}
\end{equation}

Dua hal perlu diperhatikan pada Persamaan~\eqref{eq:padding-sha}. Pertama, bit
\code{1} selalu disisipkan, bahkan bila pesan sudah berakhir tepat pada batas
blok; akibatnya pesan dengan panjang $448$ bit justru menuntut satu blok
tambahan. Aturan itu mencegah dua pesan dengan panjang berbeda menghasilkan
susunan blok yang sama. Kedua, panjang pesan \emph{ikut diproses sebagai bagian
dari data}, bukan sekadar dicatat di samping. Konsekuensinya penting: perubahan
satu bit pada panjang pesan mengubah blok akhir, sehingga menambahkan bit pada
akhir pesan tidak dapat dipakai untuk menghasilkan hash yang sama.

Setelah semua blok tersedia, fungsi kompresi diiterasi dari sebuah nilai awal
tetap (\textit{initialization vector}, IV):

\begin{equation}
h_0 = \mathrm{IV}, \qquad h_i = f(h_{i-1}, M_i) \quad \text{untuk }
i = 1, 2, \dots, t, \qquad H(m) = h_t .
\label{eq:md-iterasi}
\end{equation}

Konstruksi pada Persamaan~\eqref{eq:md-iterasi} dikenal sebagai
\textbf{konstruksi Merkle--Damgård}, dirumuskan hampir bersamaan oleh Ralph
Merkle \citep{merkle1989} dan Ivan Damg{\aa}rd \citep{damgard1989}. Kekuatannya
adalah sebuah teorema: bila fungsi kompresi $f$ tahan kolisi, maka hash $H$ yang
dibangun dengan cara itu juga tahan kolisi. Argumennya sederhana. Misalkan
ditemukan kolisi pada $H$, yaitu dua pesan dengan hasil akhir sama. Telusuri
kedua rantai kompresi dari belakang; pada suatu titik keduanya harus bertemu
dengan masukan $f$ yang berbeda tetapi keluaran yang sama, dan titik itu
memberikan kolisi pada $f$. Jadi kolisi pada fungsi kompresi menular menjadi
kolisi pada hash.

Yang perlu ditekankan, teorema itu hanya menjamin arah sebaliknya tidak berlaku:
ketahanan kolisi $H$ tidak dengan sendirinya memberi ketahanan kolisi $f$.
Karena itu penyerangan terhadap MD5 dan SHA-1 ditempuh dengan mencari kolisi
pada fungsi kompresinya, sebagaimana diperlihatkan pada
Subbagian~\ref{subsec:kriptanalisis-hash}.

\subsection{Serangan Panjang-Ekstensi}
\label{subsec:length-extension}

Konstruksi iteratif pada Persamaan~\eqref{eq:md-iterasi} meninggalkan satu cacat
struktural yang halus. Karena keadaan internal $h_t$ \emph{sama} dengan nilai
hash yang dipublikasikan, siapa pun yang mengetahui $H(m)$ sebenarnya mengetahui
seluruh keadaan mesin sesudah pesan terakhir diproses. Orang itu dapat
melanjutkan kompresi dari titik tersebut, seolah-olah telah memproses $m$
beserta padding-nya sendiri.

Akibatnya, diberikan $H(m)$ dan panjang $m$, penyerang dapat menghitung

\begin{equation}
H\big(m \,\|\, \mathrm{pad}(m) \,\|\, m'\big)
\label{eq:length-extension}
\end{equation}

untuk pesan tambahan $m'$ pilihannya sendiri, \emph{tanpa mengetahui isi $m$}.
Yang dibutuhkan hanyalah melanjutkan proses dari digest yang sudah diketahui.
Serangan ini disebut \textbf{panjang-ekstensi} (\textit{length extension}) dan
tidak menuntut perhitungan yang mahal sedikit pun: biayanya sama dengan biaya
menghitung hash pesan tambahan itu secara biasa.

Serangan panjang-ekstensi berbahaya pada skema integritas yang dibangun dengan
menggabungkan kunci dan pesan secara sederhana. Bila tag dihitung sebagai
$H(K \| m)$ dengan kunci rahasia $K$ yang dirahasiakan kedua pihak, penyerang
yang menyaksikan sebuah pesan beserta tag-nya dapat menghasilkan tag yang sah
untuk pesan yang diperpanjang --- tanpa pernah mengetahui $K$. Sisi praktis
serangan ini, lengkap dengan angka-angka nyata, diperlihatkan pada
Contoh~\ref{ex:length-extension}, sedangkan cara menutupnya dibahas pada
Subbagian~\ref{subsec:hmac}.

\subsection{Aplikasi Fungsi Hash}
\label{subsec:aplikasi-hash}

Fungsi hash muncul di banyak tempat, dan pada setiap tempat properti yang
dibutuhkan berbeda. Enam pemakaian berikut memperlihatkan rentangnya.

\begin{enumerate}
    \item \textbf{Pemeriksaan integritas berkas.} Nilai hash sebuah berkas
          diumumkan pada saluran yang tepercaya, misalnya situs resmi pengunduh.
          Pengguna yang selesai mengunduh menghitung ulang hash-nya dan
          membandingkan. Bila berbeda, berkas itu rusak atau disisipi. Properti
          yang dipakai di sini adalah second pre-image resistance: penyerang
          tidak boleh mampu menyiapkan berkas lain dengan hash yang sama.
    \item \textbf{Penyimpanan kata sandi.} Kata sandi tidak disimpan apa adanya,
          melainkan sebagai hash-nya. Bila basis data bocor, penyerang hanya
          memperoleh hash, dan untuk memulihkan kata sandi ia harus menempuh
          persoalan pre-image. Karena $2^{n}$ terlalu besar untuk ditelusuri,
          ia beralih ke serangan kamus: menghitung hash ribuan kata sandi yang
          lazim dipakai dan membandingkan hasilnya. Pertahanannya adalah
          \textbf{salt}, yaitu nilai acak yang disisipkan sebelum hashing,
          sehingga dua pengguna dengan kata sandi sama tetap menghasilkan hash
          berbeda dan serangan kamus harus diulang per pengguna.
    \item \textbf{Bukti kerja} (\textit{proof of work}). Mata uang kripto
          mensyaratkan penambang menemukan nilai acak sehingga hash sebuah blok
          berada di bawah ambang tertentu. Karena hash tidak dapat diarahkan,
          penambang harus mencoba berulang kali; biaya pencarian itulah yang
          menjadi bukti kerja.
    \item \textbf{Pohon Merkle} (\textit{Merkle tree}). Nilai hash dapat
          disusun berlapis: hash dari hash, sampai terbentuk satu akar tunggal.
          Struktur itu memungkinkan pembuktian keanggotaan satu potong data
          dengan hanya menelusuri satu lintasan, dan inilah dasar verifikasi
          ringan pada rantai blok.
    \item \textbf{Pembangkitan kunci dari kata sandi.} Fungsi hash dapat dipakai
          berulang untuk menurunkan kunci kriptografi dari kata sandi, sehingga
          menebak kata sandi menjadi mahal.
    \item \textbf{Dasar MAC dan tanda tangan digital.} Sebagaimana diperlihatkan
          pada Bagian~\ref{sec:mac} dan pada Bab~\ref{chap:tanda-tangan-pki},
          tanda
          tangan digital tidak dituliskan atas pesan utuh, melainkan atas
          hash-nya. Alasannya dua: pesan berukuran besar cukup diringkas menjadi
          $n$ bit, dan skema tanda tangan untuk pesan panjang menjadi sederhana.
          Justru karena peran itulah ketahanan kolisi menjadi syarat mutlak:
          bila penyerang dapat menyiapkan dua pesan berhash sama, tanda tangan
          untuk pesan yang tidak berbahaya dapat dipakai untuk mengesahkan pesan
          yang lain sama sekali.
\end{enumerate}

Pemakaian terakhir itu menjelaskan mengapa ketahanan kolisi, bukan pre-image
resistance, yang biasanya menentukan pilihan algoritma. Pada tanda tangan
digital, penyerang memiliki kebebasan memilih kedua pesan, sehingga biaya
$2^{n/2}$ --- bukan $2^{n}$ --- yang harus dilampauinya.

\subsection{Fungsi Hash Bukan Enkripsi}
\label{subsec:hash-bukan-enkripsi}

Karena sama-sama mengubah pesan menjadi untai bit yang tampak acak, fungsi hash
sering disalahpahami sebagai bentuk enkripsi. Tabel~\ref{tab:hash-vs-enkripsi}
memperlihatkan perbedaan yang menentukan.

\begin{table}[htbp]
    \centering
    \caption{Perbandingan fungsi hash dan enkripsi}
    \label{tab:hash-vs-enkripsi}
    \begin{tabularx}{\linewidth}{@{}l
                                   >{\raggedright\arraybackslash}X
                                   >{\raggedright\arraybackslash}X@{}}
    \toprule
    \textbf{Aspek} & \textbf{Fungsi hash} & \textbf{Enkripsi} \\
    \midrule
    Arah & satu arah; tidak ada dekripsi & dua arah; ada dekripsi \\
    \addlinespace
    Kunci & tidak memakai kunci & selalu memakai kunci \\
    \addlinespace
    Panjang keluaran & tetap, misalnya 256 bit & mengikuti panjang masukan \\
    \addlinespace
    Sifat keluaran & menyingkat; banyak masukan, satu keluaran &
    dapat dipulihkan utuh \\
    \addlinespace
    Jaminan & integritas, bukan kerahasiaan &
    kerahasiaan \\
    \addlinespace
    Bila kunci bocor & tidak berlaku: hash tidak berkunci &
    seluruh pesan terbuka \\
    \bottomrule
    \end{tabularx}
\end{table}

Dua baris terakhir Tabel~\ref{tab:hash-vs-enkripsi} menjelaskan mengapa kedua
perkakas itu tidak dapat saling menggantikan. Fungsi hash tidak menyembunyikan
apa pun: nilainya memang disiarkan terbuka, dan justru keterbukaan itu yang
membuatnya berguna sebagai pembanding. Sebaliknya enkripsi tidak memberikan
jaminan integritas: cipherteks dapat diubah oleh siapa saja yang menyadapnya,
dan perubahan itu hanya akan ketahuan sesudah didekripsi --- bila memang
ketahuan. Karena itulah pesan yang harus dirahasiakan \emph{dan} dilindungi
integritasnya memerlukan keduanya sekaligus, suatu pola yang menjadi tema
Subbagian~\ref{subsec:aead}.
